\usepackage{amsmath,amssymb,unicode-math,tikz,xcolor,sectsty,regexpatch,letltxmacro,makeidx}
%% elements of style
\definecolor{WildStrawberry}{rgb}{1,.26,.64}
\usepackage{fancyhdr}
\makeatletter
\renewcommand{\@chapapp}{Section}
\makeatother
\pagestyle{fancy}
\renewcommand{\chaptermark}[1]{\markboth{#1}{}}
\renewcommand{\sectionmark}[1]{\markright{#1}}
\fancyhead{} % clear all header fields
\fancyfoot{} % clear all footer fields
\fancyfoot[LE]{\thepage\ |\ \nouppercase{\leftmark}}
\fancyfoot[RO]{\nouppercase{\rightmark}\ |\ \thepage}
\renewcommand{\headrulewidth}{0pt}
\renewcommand{\footrulewidth}{0.4pt}
\renewcommand{\footrule}{\hbox to\headwidth{%
    \color{WildStrawberry}\leaders\hrule height \footrulewidth\hfill}}
\fancypagestyle{plain}{%
  \fancyhead{} % clear all header fields
  \fancyfoot{} % clear all footer fields
  \fancyfoot[LE]{\thepage\ |\ \nouppercase{\leftmark}}
  \fancyfoot[RO]{\nouppercase{\rightmark}\ |\ \thepage}
  \renewcommand{\headrulewidth}{0pt}
  \renewcommand{\footrulewidth}{0.4pt}
  \renewcommand{\footrule}{\hbox to\headwidth{%
      \color{WildStrawberry}\leaders\hrule height \footrulewidth\hfill}}
}
%%
\makeatletter
\newcommand{\setchapterrulecolor}[1]{\colorlet{secrulecolor}{#1}}
\setchapterrulecolor{WildStrawberry}
\xpatchcmd*{\SS@normsectionrule}% <cmd>
  {\rule}% <search>
  {\color{secrulecolor}\rule}% <replace>
  {}{}% <success><failure>
% \newcommand{\setsectionrulecolor}[1]{\colorlet{secrulecolor}{#1}}
% \setsectionrulecolor{WildStrawberry}
% \xpatchcmd*{\SS@normsectionrule}% <cmd>
%   {\rule}% <search>
%   {\color{secrulecolor}\rule}% <replace>
%   {}{}% <success><failure>
\makeatother
\chapterfont{\sectionrule{0pt}{0pt}{-5pt}{0.8pt}}
% \sectionfont{\sectionrule{0pt}{0pt}{-5pt}{0.8pt}}


%% newcommands

\DeclareMathOperator{\Dirac}{Dirac}
\DeclareMathOperator{\expit}{expit}
\DeclareMathOperator{\logit}{logit}
\DeclareMathOperator{\Rem}{Rem}
\DeclareMathOperator{\Var}{Var}

%% %% exercize symbol

\newsavebox\gearbox
\newcommand{\gearmacro}[5]{%
  % #1 number of teeths
  % #2 radius intern
  % #3 radius extern
  % #4 angle from start to end of the first arc
  % #5 angle to decale the second arc from the first 
  \foreach \i in {1,...,#1} {%
    [rotate=(\i-1)*360/#1]  (0:#2)  arc (0:#4:#2) {[rounded corners=1.5pt]
      -- (#4+#5:#3)  arc (#4+#5:360/#1-#5:#3)} --  (360/#1:#2)
  }}  
\newcommand{\makegear}{\begin{tikzpicture}
    \fill[gray] (0,0) circle (2cm);
    \draw[thick,rotate=12,fill=gray] \gearmacro{8}{2}{2.4}{20}{2};
    \draw[thick,fill=white] (0cm,0cm) circle(1.35cm);
  \end{tikzpicture}}
\sbox\gearbox{\resizebox{!}{1em}{\makegear}}
\newcommand{\gear}{\usebox{\gearbox}\;}

%% %% danger symbol

\DeclareFontFamily{U}{stixbbit}{}
\DeclareFontShape{U}{stixbbit}{m}{it}{<-> stix-mathbbit}{}
\DeclareRobustCommand{\stixdanger}{%
  {\usefont{U}{stixbbit}{m}{it}\symbol{"F6}}%
}


\newcommand{\bbO}{\mathbb{O}}
\newcommand{\bbD}{\mathbb{D}}
\newcommand{\bbP}{\mathbb{P}}
\newcommand{\bbR}{\mathbb{R}}
\newcommand{\Algo}{\widehat{\mathcal{A}}}
\newcommand{\Algora}{\widetilde{\mathcal{A}}}
\newcommand{\defq}{\doteq}
\newcommand{\calF}{\mathcal{F}}
\newcommand{\calM}{\mathcal{M}}
\newcommand{\calP}{\mathcal{P}}
\newcommand{\calO}{\mathcal{O}}
\newcommand{\calQ}{\mathcal{Q}}
\newcommand{\Exp}{\textrm{E}}
\newcommand{\IC}{\textrm{IC}}
\newcommand{\Gbar}{\bar{G}}
\newcommand{\one}{\textbf{1}}
\newcommand{\psinos}{\psi_{n}^{\textrm{os}}}
\renewcommand{\Pr}{\textrm{Pr}}
\newcommand{\Phat}{P^{\circ}}
\newcommand{\Psihat}{\widehat{\Psi}}
\newcommand{\Qbar}{\bar{Q}}
\newcommand{\tcg}[1]{\textcolor{olive}{#1}}

\allowdisplaybreaks

\makeindex

%%% Local Variables:
%%% mode: latex
%%% TeX-master: t
%%% End:
